\problema{Number of Connected Components}{323}{src/01-grafos/323-connected-components.cpp}{M}

Hay \texttt{n} nodos, numerados de \texttt{0} a \texttt{n - 1}, y una lista
\texttt{edges} de pares $[a, b]$ que representan aristas \textbf{no
dirigidas}. Queremos contar cuántos componentes conexos tiene el grafo.

\paragraph{Intuición.} Empezamos suponiendo lo peor: que cada nodo vive
solo, aislado, formando su propio componente. Eso nos da \texttt{n}
componentes de entrada. Cada arista es una oportunidad de fusión: si conecta
dos nodos que ya viven en el mismo componente, no cambia nada (la arista es
``redundante'' para efectos de contar componentes); pero si conecta dos
nodos de componentes distintos, esos dos componentes se convierten en uno
solo y el contador baja en $1$. Al final de procesar todas las aristas, el
contador tiene la respuesta.

\begin{center}
\ttfamily
\begin{tabular}{c}
$n=5$, edges $=[[0,1],[1,2],[3,4]]$\\[4pt]
0 --- 1 --- 2 \qquad\qquad 3 --- 4\\[4pt]
2 componentes: $\{0,1,2\}$ y $\{3,4\}$
\end{tabular}
\end{center}

\paragraph{Solución (Union-Find / DSU).} La estructura ideal para este
``ir fusionando conjuntos'' es \textbf{Union-Find} (Disjoint Set Union):
cada nodo apunta a un padre, y la \emph{raíz} de un nodo identifica a qué
componente pertenece. Dos operaciones:
\begin{itemize}
\item \texttt{find(x)}: sigue los punteros de padre hasta la raíz del
  conjunto de \texttt{x}. Aplicamos \emph{compresión de caminos}: cada nodo
  visitado en el camino se re-apunta directamente a la raíz, para que
  futuras búsquedas sean casi instantáneas.
\item \texttt{unir(a, b)}: si \texttt{find(a)} y \texttt{find(b)} son
  distintas, cuelga una raíz de la otra usando \emph{unión por rango}
  (la raíz de árbol más pequeño cuelga de la más grande, para mantener los
  árboles bajos) y decrementa el contador de componentes.
\end{itemize}
Inicializamos cada nodo como su propio padre (\texttt{n} componentes) y
recorremos las aristas una vez, uniendo cuando sea necesario.

\lstinputlisting{src/01-grafos/323-connected-components.cpp}

\complejidad{$O(n + e\cdot\alpha(n))$}{$O(n)$}

\noindent donde $\alpha(n)$ es la función inversa de Ackermann (crece tan
lento que para cualquier $n$ práctico es prácticamente una constante), y
$e$ es el número de aristas.

\paragraph{Notas de entrevista (Amazon).} Este problema es el ejemplo
canónico de Union-Find y es hermano directo del \textbf{261. Graph Valid
Tree}: ahí, en vez de contar componentes, se pregunta si el grafo termina
siendo un único árbol (exactamente $n-1$ aristas y ninguna arista que una
dos nodos ya en el mismo componente, es decir, ningún ciclo). También es
primo del patrón de ``contar islas'' (200, 695): en lugar de recorrer una
cuadrícula con DFS/BFS marcando visitados, aquí el grafo viene dado
explícitamente como lista de aristas, pero la idea de fondo es la misma:
contar cuántos grupos conectados existen. De hecho, este problema también
se puede resolver con DFS o BFS clásico, construyendo una lista de
adyacencia y contando cuántas veces arrancamos una exploración desde un
nodo no visitado. Menciona ambos enfoques en la entrevista: DFS/BFS es más
directo de explicar, pero Union-Find es más eficiente cuando las aristas
llegan de forma incremental (por ejemplo, si el entrevistador pregunta un
seguimiento tipo ``¿y si las aristas se agregan una a una y quieres saber
el número de componentes después de cada una?''). Casos borde: sin
aristas (cada nodo es su propio componente), un solo nodo, y aristas
repetidas o que no cambian el número de componentes.
